\documentclass[conference]{IEEEtran}
\IEEEoverridecommandlockouts

% Configuración para español
\usepackage[spanish,es-nodecimaldot,es-noshorthands]{babel}
\usepackage[utf8]{inputenc}
\usepackage[T1]{fontenc}

% Paquetes estándar IEEE
\usepackage{cite}
\usepackage{amsmath,amssymb,amsfonts}
\usepackage{algorithmic}
\usepackage{graphicx}
\usepackage{textcomp}
\usepackage{xcolor}
\usepackage{float}
\usepackage{listings}
\usepackage{multirow}
\usepackage[hidelinks]{hyperref} % hidelinks elimina las cajas de color en los enlaces

% ---------- Configuración de listados de código ----------
\definecolor{codegray}{rgb}{0.45,0.45,0.45}
\definecolor{codeblue}{rgb}{0.10,0.10,0.60}
\lstset{
	basicstyle=\ttfamily\scriptsize,
	keywordstyle=\color{codeblue}\bfseries,
	commentstyle=\color{codegray}\itshape,
	stringstyle=\color{codegray},
	numbers=left,
	numberstyle=\tiny\color{codegray},
	numbersep=5pt,
	frame=single,
	breaklines=true,
	breakatwhitespace=true,
	showstringspaces=false,
	tabsize=2,
	captionpos=b,
	literate=%
	{á}{{\'a}}1 {é}{{\'e}}1 {í}{{\'i}}1 {ó}{{\'o}}1 {ú}{{\'u}}1
	{Á}{{\'A}}1 {É}{{\'E}}1 {Í}{{\'I}}1 {Ó}{{\'O}}1 {Ú}{{\'U}}1
	{ñ}{{\~n}}1 {Ñ}{{\~N}}1 {¿}{{?`}}1 {¡}{{!`}}1
}

% ---------- Espacio reservado para imágenes faltantes ----------
\newcommand{\imgph}[2]{%
	\fbox{\begin{minipage}[c][#2][c]{0.92\columnwidth}
			\centering\textcolor{codegray}{\footnotesize [Espacio para imagen: #1]}
	\end{minipage}}%
}

\def\BibTeX{{\rm B\kern-.05em{\sc i\kern-.025em b}\kern-.08em
	T\kern-.1667em\lower.7ex\hbox{E}\kern-.125emX}}

\begin{document}

\title{Propuesta de Proyecto Final: \\
	\large [Nombre del Proyecto] \\
	Sistemas Embebidos - Avance 1}

\author{
	\IEEEauthorblockN{Belmont Muñoz Samuel}
	\IEEEauthorblockA{\textit{Ingeniería en Computación} \\
		\textit{Facultad de Ingeniería}\\
		No. Cuenta: 317318289}
	\and
	\IEEEauthorblockN{Domínguez Palacios Jesús Alejandro}
	\IEEEauthorblockA{\textit{Ingeniería en Computación} \\
		\textit{Facultad de Ingeniería}\\
		No. Cuenta: 421066205}
	\and
	\IEEEauthorblockN{Dorantes Mejía Emmanuel}
	\IEEEauthorblockA{\textit{Ingeniería en Computación} \\
		\textit{Facultad de Ingeniería}\\
		No. Cuenta: 313004766}
	\and
	\IEEEauthorblockN{García Olvera Emiliano}
	\IEEEauthorblockA{\textit{Ingeniería en Computación} \\
		\textit{Facultad de Ingeniería}\\
		No. Cuenta: 314256009}
	\and
	\IEEEauthorblockN{Hernández Alonso Antonio Adonahí}
	\IEEEauthorblockA{\textit{Ingeniería en Computación} \\
		\textit{Facultad de Ingeniería}\\
		No. Cuenta: 320257780}
	\and
	\IEEEauthorblockN{Jiménez Reyes Alan Yael}
	\IEEEauthorblockA{\textit{Ingeniería en Computación} \\
		\textit{Facultad de Ingeniería}\\
		No. Cuenta: 320026726}
}

\maketitle

\begin{center}
	\vspace{-0.3cm}
	{\small
		\textbf{Grupo:} 6  | 
		\textbf{Fecha de entrega:} 19/10/2026  | 
		\textbf{Avance:} 1 de 4
	}
	\vspace{0.2cm}
\end{center}

\begin{abstract}
	Este documento presenta la propuesta de diseño para [nombre del proyecto], un sistema embebido desarrollado como proyecto final de la asignatura de Sistemas Embebidos. Se describe el problema a resolver, los usuarios objetivo, el alcance del sistema, los requisitos funcionales y no funcionales con criterios de aceptación numéricos, la arquitectura del sistema mediante diagrama de bloques, la justificación de la plataforma de hardware y software seleccionada, la bitácora de decisiones de diseño, la lista de materiales con análisis de riesgos, y el plan de trabajo con asignación de responsabilidades. El sistema está diseñado para [breve descripción del propósito].
\end{abstract}

\begin{IEEEkeywords}
	Sistemas embebidos, Raspberry Pi, [otros keywords del proyecto]
\end{IEEEkeywords}

\section{Problema, Usuario y Alcance}

\subsection{Descripción del Problema}

[Describir el problema que resuelve el sistema. ¿Qué necesidad o situación motiva este proyecto? ¿Por qué es importante resolverlo?]

\subsection{Usuario Objetivo}

[¿Quién utilizará este sistema? Describir el perfil del usuario, sus necesidades específicas y el contexto de uso.]

\subsection{Alcance del Proyecto}

\subsubsection{Dentro del Alcance}

El sistema incluirá las siguientes funcionalidades:

\begin{itemize}
	\item {[Funcionalidad 1]}
	\item {[Funcionalidad 2]}
	\item {[Funcionalidad 3]}
	\item {[Funcionalidad 4]}
	\item {[Funcionalidad 5]}
\end{itemize}

\subsubsection{Fuera del Alcance}

Las siguientes funcionalidades NO serán implementadas en este proyecto:

\begin{itemize}
	\item {[Exclusión 1]}
	\item {[Exclusión 2]}
	\item {[Exclusión 3]}
\end{itemize}

\section{Requisitos del Sistema}

Esta sección presenta los requisitos funcionales y no funcionales del sistema, cada uno con su criterio de aceptación numérico medible.

\subsection{Requisitos Funcionales}

\begin{enumerate}
	\item \textbf{[Nombre del Requisito 1]} \\
	      Descripción: [Descripción del requisito] \\
	      \textit{Criterio de aceptación:} [Criterio numérico medible, ej: "El sistema debe detectar eventos en menos de 100ms"]

	\item \textbf{[Nombre del Requisito 2]} \\
	      Descripción: [Descripción del requisito] \\
	      \textit{Criterio de aceptación:} [Criterio numérico medible]

	\item \textbf{[Nombre del Requisito 3]} \\
	      Descripción: [Descripción del requisito] \\
	      \textit{Criterio de aceptación:} [Criterio numérico medible]
\end{enumerate}

\subsection{Requisitos No Funcionales}

\begin{enumerate}
	\item \textbf{[Nombre del Requisito 4]} \\
	      Descripción: [Descripción del requisito] \\
	      \textit{Criterio de aceptación:} [Criterio numérico medible, ej: "Consumo de energía menor a 5W en operación continua"]

	\item \textbf{[Nombre del Requisito 5]} \\
	      Descripción: [Descripción del requisito] \\
	      \textit{Criterio de aceptación:} [Criterio numérico medible]
\end{enumerate}

\section{Diagrama de Bloques}

El diagrama debe mostrar claramente: sensores, actuadores, procesadores, buses de comunicación (GPIO, I2C, SPI, UART) y el flujo de datos entre componentes.

\begin{figure}[H]
	\centering
	\imgph{Diagrama de bloques del sistema}{6cm}
	\caption{Arquitectura del sistema con sensores, actuadores, procesadores, buses y flujo de datos.}
	\label{fig:diagrama_bloques}
\end{figure}

\textbf{Descripción:} {[Breve descripción textual del diagrama explicando los componentes principales, sus conexiones y el flujo de datos]}

\section{Plataforma y Sistema Operativo}

\subsection{Plataforma de Hardware}

\textbf{Plataforma elegida:} {[Raspberry Pi 4 / STM32 / Ambas]}

\textbf{Justificación:}

\begin{itemize}
	\item \textbf{Determinismo}: {[Análisis de requisitos de tiempo real y cómo la plataforma los cumple]}
	\item \textbf{Huella de memoria}: {[Cifras de RAM y almacenamiento necesarios vs disponibles]}
	\item \textbf{Multitarea}: {[Cantidad de tareas concurrentes y cómo la plataforma las soporta]}
\end{itemize}

\subsection{Categoría de Sistema Operativo}

\textbf{SO seleccionado:} {[Linux embebido / RTOS / Bare-metal]}

\textbf{Justificación:}

\begin{itemize}
	\item \textbf{Determinismo}: {[¿Requiere garantías de tiempo real? Cifras de latencia máxima]}
	\item \textbf{Huella de memoria}: {[Overhead del SO en KB/MB vs memoria disponible]}
	\item \textbf{Multitarea}: {[Tipo de planificación requerida y soporte del SO]}
\end{itemize}

\section{Bitácora de Decisiones de Diseño}

Esta sección documenta las decisiones técnicas tomadas durante las clases 7 a 12, aplicadas específicamente a nuestro proyecto. Cada entrada incluye cifras y justificaciones cuantitativas.

\subsection{Clase 7: Kickoff de proyecto + Panorama de SO embebidos}

\textbf{Pregunta:} ¿Qué categoría de SO embebido usarían para su proyecto — bare-metal, RTOS o Linux embebido — y por qué?

\textbf{Respuesta aplicada al proyecto:} \\
{[Respuesta específica con cifras y detalles técnicos relacionados con el proyecto]}

\textbf{Autor:} {[Nombre del estudiante]} \\
\textbf{Fecha:} 09/09/2026

\subsection{Clase 8: Planificación (scheduling)}

\textbf{Pregunta:} ¿Qué tareas tendría su proyecto, y qué prioridad relativa le darían a cada una? ¿Alguna necesita planificación preemptiva sí o sí?

\textbf{Respuesta aplicada al proyecto:} \\
{[Respuesta específica con cifras y detalles técnicos relacionados con el proyecto]}

\textbf{Autor:} {[Nombre del estudiante]} \\
\textbf{Fecha:} 10/09/2026

\subsection{Clase 9: Sincronización}

\textbf{Pregunta:} ¿Qué recurso compartido tendría su proyecto entre dos tareas, y qué mecanismo de sincronización usarían para protegerlo?

\textbf{Respuesta aplicada al proyecto:} \\
{[Respuesta específica con cifras y detalles técnicos relacionados con el proyecto]}

\textbf{Autor:} {[Nombre del estudiante]} \\
\textbf{Fecha:} 17/09/2026

\subsection{Clase 10: Gestión de Memoria en SO embebidos}

\textbf{Pregunta:} ¿Su proyecto se parece más al Caso A (asignaciones frecuentes y predecibles) o al Caso B (una asignación grande y rara)? ¿Qué estrategia de memoria usarían, con qué números?

\textbf{Respuesta aplicada al proyecto:} \\
{[Respuesta específica con cifras y detalles técnicos relacionados con el proyecto]}

\textbf{Autor:} {[Nombre del estudiante]} \\
\textbf{Fecha:} 23/09/2026

\subsection{Clase 11: Comunicación entre Procesos (IPC)}

\textbf{Pregunta:} ¿Su proyecto tendría más de un proceso comunicándose? Si es así, ¿qué mecanismo de IPC usarían y por qué?

\textbf{Respuesta aplicada al proyecto:} \\
{[Respuesta específica con cifras y detalles técnicos relacionados con el proyecto]}

\textbf{Autor:} {[Nombre del estudiante]} \\
\textbf{Fecha:} 24/09/2026

\subsection{Clase 12: Kernel, syscalls y drivers}

\textbf{Pregunta:} ¿Qué tareas de su proyecto tienen una temporización que un proceso de usuario no puede garantizar? Para cada una, ¿usarían una syscall por evento, mmap de registros o un periférico de hardware configurado por su driver?

\textbf{Respuesta aplicada al proyecto:} \\
{[Respuesta específica con cifras y detalles técnicos relacionados con el proyecto]}

\textbf{Autor:} {[Nombre del estudiante]} \\
\textbf{Fecha:} 07/10/2026

\section{Lista de Materiales y Análisis de Riesgos}

\subsection{Bill of Materials (BOM)}

La Tabla \ref{tab:bom} presenta la lista completa de materiales necesarios para el proyecto, con información de disponibilidad confirmada.

\begin{table}[H]
	\centering
	\caption{Lista de materiales del proyecto}
	\label{tab:bom}
	\small
	\begin{tabular}{|p{3cm}|c|c|p{2.3cm}|}
		\hline
		\textbf{Componente} & \textbf{Cant.} & \textbf{Precio} & \textbf{Disponibilidad} \\
		\hline
		Raspberry Pi 4 (4GB) & 1 & \$X & {[Confirmada/ Disponible/ A comprar]} \\
		\hline
		{[Sensor 1]} & X & \$X & {[Estado]} \\
		\hline
		{[Sensor 2]} & X & \$X & {[Estado]} \\
		\hline
		{[Actuador 1]} & X & \$X & {[Estado]} \\
		\hline
		{[Componente adicional]} & X & \$X & {[Estado]} \\
		\hline
		Cables y conectores & - & \$X & Disponible \\
		\hline
		\textbf{Total estimado} & & \textbf{\$XXX} & \\
		\hline
	\end{tabular}
\end{table}

\subsection{Análisis de Riesgos y Mitigación}

La Tabla \ref{tab:riesgos} identifica los principales riesgos del proyecto y las estrategias de mitigación planificadas.

\begin{table}[H]
	\centering
	\caption{Riesgos identificados y plan de mitigación}
	\label{tab:riesgos}
	\small
	\begin{tabular}{|p{2.3cm}|p{2.3cm}|p{2.7cm}|}
		\hline
		\textbf{Riesgo} & \textbf{Impacto} & \textbf{Mitigación} \\
		\hline
		{[Riesgo 1: ej. Retraso en llegada de componentes]} & {[Descripción del impacto]} & {[Plan de mitigación específico]} \\
		\hline
		{[Riesgo 2: ej. Falla de comunicación I2C]} & {[Descripción del impacto]} & {[Plan de mitigación específico]} \\
		\hline
		{[Riesgo 3: ej. Insuficiente capacidad de procesamiento]} & {[Descripción del impacto]} & {[Plan de mitigación específico]} \\
		\hline
		{[Riesgo 4]} & {[Descripción del impacto]} & {[Plan de mitigación específico]} \\
		\hline
	\end{tabular}
\end{table}

\section{Plan de Trabajo}

\subsection{Cronograma de Hitos}

El proyecto se desarrollará en 4 avances principales:

\begin{itemize}
	\item \textbf{Avance 1 (19/oct/2026)}: Propuesta + diagrama de bloques
	\item \textbf{Avance 2 (04/nov/2026)}: Avance de al menos 1 módulo funcional
	\item \textbf{Avance 3 (09/nov/2026)}: Avance de integración completa
	\item \textbf{Avance 4 (23/nov/2026)}: Demostración final del prototipo + reporte técnico
\end{itemize}

\subsection{Distribución de Tareas}

La Tabla \ref{tab:tareas} muestra la asignación de responsabilidades por tarea y por avance.

\begin{table}[H]
	\centering
	\caption{Distribución de tareas y responsables}
	\label{tab:tareas}
	\small
	\begin{tabular}{|p{1.6cm}|p{3.2cm}|p{2.2cm}|}
		\hline
		\textbf{Avance} & \textbf{Tarea} & \textbf{Responsable} \\
		\hline
		\multirow{3}{*}{Avance 2} & {[Tarea 1]} & {[Nombre]} \\
		\cline{2-3}
		& {[Tarea 2]} & {[Nombre]} \\
		\cline{2-3}
		& {[Tarea 3]} & {[Nombre]} \\
		\hline
		\multirow{3}{*}{Avance 3} & {[Tarea 1]} & {[Nombre]} \\
		\cline{2-3}
		& {[Tarea 2]} & {[Nombre]} \\
		\cline{2-3}
		& {[Tarea 3]} & {[Nombre]} \\
		\hline
		\multirow{3}{*}{Avance 4} & {[Tarea 1]} & {[Nombre]} \\
		\cline{2-3}
		& {[Tarea 2]} & {[Nombre]} \\
		\cline{2-3}
		& Integración final & Todos \\
		\hline
	\end{tabular}
\end{table}

\section{Declaración de Uso de Inteligencia Artificial}

Durante el desarrollo de este proyecto se utilizaron las siguientes herramientas de IA:

\begin{itemize}
	\item \textbf{Claude Code (Anthropic)}: Utilizado para [describir para qué se usó: ej. asistencia en estructura de código, revisión de sintaxis, generación de documentación, etc.]
	\item \textbf{[Otra herramienta si aplica]}: [Descripción de uso]
\end{itemize}

Todas las decisiones de diseño, arquitectura y justificaciones técnicas fueron realizadas por el equipo. Las herramientas de IA se utilizaron como asistentes para [redacción/código/documentación], pero el contenido técnico y conceptual es autoría del equipo.

\begin{thebibliography}{00}
	\bibitem{rpi_docs} Raspberry Pi Foundation. (2024). ``Raspberry Pi Documentation''. Disponible en: https://www.raspberrypi.com/documentation/
	\bibitem{linux_embedded} Yaghmour, K. et al. (2008). \textit{Building Embedded Linux Systems}, 2nd ed. O'Reilly Media.
	\bibitem{embedded_systems} Lee, E. A., Seshia, S. A. (2017). \textit{Introduction to Embedded Systems: A Cyber-Physical Systems Approach}, 2nd ed. MIT Press.
	\bibitem{[referencia_especifica]} [Autor]. ([Año]). ``[Título]''. [Fuente].
\end{thebibliography}

\vspace{0.5cm}

\noindent
\textbf{Enlace al repositorio del proyecto:} \\
\url{https://github.com/[usuario]/[repositorio]}

\end{document}
